\honest{Excluding the Supernatural}{Eliezer Yudkowsky}{2008}

Occasionally, I say, you hear that creationism is excluded from science a priori because it invokes the ``supernatural.'' \nb{A federal court said so in \textsc{Kitzmiller v. Dover} (2005).} So creationism could be true and still be barred? ``It seems clear enough,'' I say, that the idea comes from a wish to avoid a confrontation between science and religion, so that no one has to say religious claims were tested and found false. \nb{The post gives no evidence for this motive. There is also a legal reason: the United States Constitution bars teaching religion presented as science, not false science. And the Dover ruling did say: ``ID's negative attacks on evolution have been refuted by the scientific community.''} Excluding Intelligent Design a priori would justify the creationists' complaint. \nb{Larry Laudan made a similar objection to the Arkansas creation-science ruling of 1982.}

What exactly do they mean by ``supernatural''? I do not report their answer. \nb{Richard Carrier reports that in science and law the word is used for ``the untestable,'' and argues against that use. On that meaning the rule excludes untestable claims and says nothing about whether they are true.} The best definition I have heard is Carrier's: a supernatural explanation appeals to ``ontologically basic mental things.'' \nb{``Uncritical Supercriticality'' used this definition earlier without naming Carrier.}

But then, it seems, we have ``very good'' grounds for excluding the supernatural a priori. ``My thesis is that non-reductionism is a confusion.'' If a 747 existed apart from its quarks, what would you observe? \nb{Philosophers of emergence give an answer: new forces acting only in certain configurations. Brian McLaughlin judged in 1992 that there is ``not a scintilla of evidence'' for them, a verdict from observation.} A brain made of quarks, I add, can only predict what quarks can, so ``You can never get Bayesian confirmation for a hypothesis of irreducibility.'' \nb{The same day, in ``Psychic Powers,'' Yudkowsky conceded Benja Fallenstein's objection to this sentence.}

Should ID then be excluded a priori? ``Of course not.'' For every irreducible God there is a reducible alien with the same predictions, and the predictions can be tested, at the Red Sea and in biology. The Jehovah model is ``coherent, falsifiable, and wrong,'' excluded a posteriori. Natural selection, I note, ``was not hailed as the long-awaited Creator,'' because ``It wasn't fundamentally mental.'' \nb{In ``An Alien God,'' Yudkowsky had said that true believers in a vague deity would have recognized evolution as their creator. This is the reason they could have given.}

Reductionism is ``just disbelief in fundamentally complicated things.'' If that sounds like an oxymoron, that is why I think non-reductionism is a confusion. Still, ``the ultimate rule of science is to look and see.''

My corollary: a designer of artificial general intelligence who speaks of religious beliefs ``in respectful tones'' knows almost nothing of reducing the mental, unless they are an idiot savant or ``outright lying.'' \nb{No case is given of a designer who fits the corollary.}
